% Reusable results and decision criteria; derivations are in publication paper IX.
\section{重み三のPoisson空間：境界と留数の再利用}
\label{sec:poisson-research-notes-v32}

\subsection{全方程式と、商を取らないという規約}
\[
 c\in S_3,\qquad D\in\mathfrak{sl}_3,\qquad
 [F]\in S_4^{\oplus3}/S_3x,
 \qquad S_d=\mathbb C[x,y,z]_d
\]
を座標とし、
\[
 D(c)=0,\qquad
 M_{c,D}(F)=2F(c)-(x\times D)\cdot F-c\operatorname{div}F=0
\]
を課す。全ての分類は、この五十三変数・三十八二次式のスキームから出発する。
三十八は最小二次生成元の個数であり、余次元ではない。
以下の成分次元はスカラー倍や重み付き自己同型で商する前の値である。
基底 $D(c)=0$ が非被約なら、その冪零方向を混合写像に残したまま計算する必要がある。

再利用できる幾何的辞書は
\[
 \alpha=c\,dw-w\,dc+w(x\times D)\cdot dx+(F\times x)\cdot dx
\]
である。これは重み六の水平一形式で、Poisson方程式は
$\alpha\wedge d\alpha=0$ とスキーム論的に同値である。
三次式による $w$ の移動は、混合核に実際の核方向を作る。
固定した一次方向の量子化と、指定済みVK族の高次係数の補完とは区別する。

\subsection{認定済み成分の表}
\begin{center}\small
\begin{tabular}{lrl}\toprule
成分&次元&再利用する特徴\\\midrule
$\mathbb P^2$ からの引戻し&35&一般点で $c=D=0$\\
$R(1,5)$&28&直線・二次曲線の共鳴比 $a=2$\\
$L(1,1,4)$&23&対数三角形の一つの共鳴壁\\
$R(2,4)$&22&直線・二次曲線の共鳴比 $a=-1$\\
$R(3,3)$&19&一般の $c$ が滑らか、$D=0$\\
$L(1,2,3)$&19&一般の直線・二次曲線型\\
$L(1,1,1,3)$&19&一般の非共点三直線型\\
$\mathbb P(1,1,3)$ からの引戻し&19&一般点の横断環の長さ六\\
接線・二次曲線型 $a=-5/2$&18&一般点で滑らか、有理第一積分\\
接線・二次曲線型 $a=-3/2$&18&一般点の長さ二、量子中心は $D_7$ 型\\\bottomrule
\end{tabular}
\end{center}
十成分のうち八成分は一般点で滑らかであり、残る二成分の長さは六と二である。
十個が全てであること、三十五が最大次元であることは未証明である。
$R,L$ は古典的な有理型・対数型の構成を表し、構成自体を新規性とはしない。
本研究で必要なのは、この重みのJacobiスキームでの成分性・厚み・交叉の認定である。
証明と引用の位置付けは公表稿IX\cite{PaperIX}にまとめる。

\subsection{共点三直線の重み三での数値}
節~\ref{sec:allm-notes-v33}の全重みの結果を $m=3$ に特殊化すると、
\[
 \mathbb C[[s_1,\ldots,s_{18},a,b,u]]/(ab,b^2,a^5u)
\]
となる。一般点の長さは六、primary成分の交叉の長さは五、一般切断の長さは七である。
元の有理数証明書も比較用に保持する。
$(A,B)=(2,5)$ で、自由な二変数四次係数を $h_{ij}=[x^iy^j]F_z$ とすると、
Schur消去の非零式は
\[
 a^5L,\ a^6L,\qquad
 L=(20h_{13}-100h_{22}+300h_{31}-700h_{40})/9
\]
であり、$F=(0,0,x^4)$ の係数は $-700/9$ である。
この式は全重み検証器が従来の規約と一致することを確認する基準になる。
全重みの証明ではこの特定の係数値を必要とせず、重量と三交点評価によって
先頭係数の非零性を保証する。

\subsection{非共点三直線：五つの局所枝}
$c=3xyz,D=F=0$ の完備局所環は
\[
 \mathbb C[[z_1,\ldots,z_{16}]]\widehat\otimes
 \frac{\mathbb C[[a_i,\delta_i,\rho_i\ (1\le i\le3)]]}
 {(\sum\delta_i,\sum\rho_i,a_i\delta_i,a_i^2\rho_i,
     \delta_1\rho_2-\delta_2\rho_1)}.
\]
ここでは $a_1\rho_1$ が非零の平方零元である。
被約化の局所枝は、$R(3,3)$、$L(1,1,1,3)$、および
$L(1,2,3)$ の三つの局所枝で、全て十九次元となる。
後者を別々の大域成分と数えない。
交叉次元は、三次曲線型と対数型が十六、三次曲線型と各直線・二次曲線枝が十七、
対数型と各直線・二次曲線枝が十八、異なる直線・二次曲線枝どうしが十六である。
残余二方向は既知の対数型の境界に収まり、この点だけから新たな大域成分は生じない。

\subsection{被約な共鳴交叉と明示量子化}
\[
 c=x(x^2+yz),\qquad D=a(y\partial_y-z\partial_z)
\]
では混合核次元が、一般の $a$ で十、$a=2$ で二十、$a=-1$ で十四となる。
局所環は滑らかな十八変数を除いて
\[
 \mathbb C[[t,v_1,\ldots,v_r]]/(tv_1,\ldots,tv_r),
 \qquad r=10\ (a=2),\quad r=4\ (a=-1)
\]
である。一般対数三角形の壁 $3\alpha+t=0$ では同じ形の $r=5$ が現れる。
これらは全て被約であり、共点三直線の厚みとは異なる局所模型である。

$r=5$ の一般点は、二変数四次式 $f$ と二つの非零係数 $a,b$ による
\[
 S=\frac{\mathbb C\langle y,z,x,w\rangle}
 {(yz-zy,\ xy-ayx,\ xz-bzx,\ wy-a^{-1}yw,\ wz-b^{-1}zw,\ wx-abxw+f)}
\]
で量子化できる。これは重み $(1,1,1,3)$ のPBW基底を持つnoetherian整域で、
AS正則次元四、通常の意味で4-Calabi--Yau、Gorenstein parameter六である。
補助次数 $\deg' x=1,\deg' w=-1,\deg' y=\deg' z=0$ は、
Ore定理で未決のNakayama平行移動項を消すために有効である。
Orlovの定理により標準六項列を持ち、$S_{(0,1,3,4)}$ の圏は実際の半直交成分となる。
任意の形式量子化のAS正則性や、全てのPoisson成分の明示量子化はまだ示していない。

\subsection{$D=0$ の全混合核}
非零三次式 $c=\prod f_i^{m_i}$ に対し
\[
 c_{\mathrm{red}}=\prod f_i,\quad d=\deg c_{\mathrm{red}},\quad
 h=c^2/c_{\mathrm{red}}
\]
と置く。Euler商の全解は
\[
 \frac{(F\times x)\cdot dx}{c^2}
 =d\left(\frac ph\right)+\sum_i\lambda_i\frac{df_i}{f_i},
 \quad p\in S_{6-d}/\mathbb Ch,\quad
 \sum_i\deg(f_i)\lambda_i=0
\]
で与えられる。核次元は $\dim S_{6-d}+r-2$、三次移動の像は常に九次元である。
既約三次、直線と二次曲線、相異なる三直線、二重直線と別の直線、三重直線の順に
核次元は $9,10,11,15,20$ となる。
ただしこのファイバーの次元増加だけでは、全Poisson空間の成分性は分からない。
$D$ を含む横断変形と混合方程式を同時に調べる必要がある。

\subsection{標識と自己同型を保持するための注意}

三次式による $w$ の変更で混合項を消す議論は、重み付き座標変換を許した標準形である。
その座標変換が特殊ファイバー上で恒等的とは限らないので、
標識付き変形の同値と取り違えてはいけない。
また、二次 Poisson 錐と外側の miniversal hull の同型は、
普遍量子化と hull の座標の選択に依存する。
指定された VK 族の補完条件をその座標で書く際には、その族を表す高次係数も追跡する必要がある。

\subsection{別の対象へ一般化する際の量子化経路}

今回の重み付き射影スタックでは、BV 恒等式
$\delta(E\wedge U)+E\wedge\delta U=6U$ によるアフィン持ち上げが最短の証明である。
Euler 場を持たない対象への一般化には、Pridham の量子化定理
\cite[Corollary~3.38]{Pridham2025} が別の入力になる。
同定理は完全な余接複体を持つ導来 Artin スタックの Poisson 構造を量子化し、
強準コンパクトの場合には perfect 圏の curved $A_\infty$ 変形を与える。

強例外列がある場合、各対象の曲率を逐次除く障害は自己 $\operatorname{Ext}^2$ に入る。
これが消滅すれば対象を持ち上げられ、全 morphism cohomology が零次に集中する場合には
最小 $A_\infty$ 構造に $m_2$ 以外の演算が残らない。
この経路は一般化に有用だが、本稿の weighted lift の証明の代わりに
曖昧な「スタック上の formality」を引用する必要はない。

\subsection{変形関手の比較をさらに強める場合}

発散零の重み零多重ベクトル場は Schouten 括弧で閉じ、各次数で
$H^0(X,\wedge^pT_X)$ と同一視できる。
アフィン Kontsevich $L_\infty$ 写像をこの部分代数に制限し、
得られる重みを保つ polydifferential cochain を六対象の可合成列に評価すれば、
外側の Hochschild 複体への $L_\infty$ 写像を構成できる。
これを用いて変形groupoidそのものの比較を述べるには、
評価の第一成分と untwisted HKR の自然性、零次 cochain から生じる内部自己同型、
および負次数を扱う規約を明記する必要がある。
現時点の hull 同定には、この追加のgroupoidレベルの主張は不要である。

